\documentclass[11pt]{article}

\usepackage[margin=1in]{geometry}
\usepackage[T1]{fontenc}
\usepackage[utf8]{inputenc}
\usepackage{lmodern}
\usepackage{amsmath,amssymb}
\usepackage{hyperref}

\setlength{\parindent}{0pt}
\setlength{\parskip}{6pt}

\title{Stream Cipher Blind Time Long Code (SC-BLTC) Protocol Specification}
\date{}

\begin{document}
\maketitle
\tableofcontents

\section{System Parameters}

The following global parameters are predefined:

\begin{itemize}
\item $F_s$ (Sampling Rate): $25\,\text{kHz}$. This is the time-domain resolution of the digital baseband waveform.
\item $R_c$ (Chip Rate): $5\,\text{kcps}$. This results in a signal bandwidth of approximately $5$--$6\,\text{kHz}$ (depending on the RRC roll-off factor $\alpha=0.25$). $T_c = 1/R_c$.
\item $\mathit{OSF}$ (Oversampling Factor): $\mathit{OSF} = F_s / R_c = 5$. $\mathit{OSF}$ is fixed as an integer.
\item $\mathit{SF}$ (Spreading Factor): $1024$. The processing gain per spread symbol is $G_p = 10\log_{10}(1024) \approx 30\,\text{dB}$.
\item $k$ (Bits per Walsh Symbol): Fixed at $k=8$.
\item $M_W$ (Walsh Orthogonal Set Size): $M_W=2^k=256$ (256-ary orthogonal modulation).
\item $R_{\text{sym}}$ (Spread Symbol Rate): $R_{\text{sym}}=R_c/\mathit{SF}=5000/1024\approx 4.88\,\text{sym/s}$.
\item $R_{\text{bit}}$ (Post-spreading Bit Rate): $R_{\text{bit}}=k\cdot R_{\text{sym}}\approx 39.06\,\text{bps}$ (excluding preamble/pilot symbols and FEC overhead).
\item $R_{\text{FEC}}$ (Polar Code Rate): Fixed at $1/2$. The information rate is $R_{\text{info}} = R_{\text{bit}} \times R_{\text{FEC}}$ (excluding Header/CRC/padding overhead and preamble/pilot symbol overhead).
\item $K_{\text{sec}}$ (Shared Secret Key): A 256-bit key, pre-shared between the transmitter and receiver.
\item $T_{\text{epoch}}$ (Time Epoch): Defined as the Unix Epoch (1970-01-01 00:00:00 UTC).
\item $\mathit{IV}_{\text{res}}$ (Time Counter Resolution): Defines the time granularity for generating the spreading code seed, fixed at $1\,\text{ms}$.
\item \textbf{TimeIndex Search Window}: The TimeIndex hypothesis set is
  \[
    \mathcal{T}=\{\mathit{TI}_{\min},\mathit{TI}_{\min}+1,\dots,\mathit{TI}_{\min}+N_{\mathit{TI}}-1\}.
  \]
  Protocol operation requires $\mathit{TI}_{\text{tx}}\in\mathcal{T}$.
\item $M$ (FEC Codeword Length, coded bits): Fixed value $M=512$.
\item $K$ (FEC Information Bit Length, uncoded bits): Fixed value $K=256$, satisfying $K=M\cdot R_{\text{FEC}}$.
\item \textbf{FEC}: Uses an Arikan Polar code with $(N,K)=(512,256)$ and CRC-aided SCL (CA-SCL) decoding.
\item \textbf{CRC}: A CRC-32C is appended inside $U$ and is used both for post-decoding verification and for CA-SCL list selection.
\item \textbf{Number of Preamble Symbols}: Fixed at $N_{\text{pre}}=2$.
\item \textbf{Number of Data Symbols}: $N_{\text{data}}=M/k=64$.
\item \textbf{Number of Pilot Symbols}: Fixed at $N_{\text{pilot}}=16$. Pilot symbols use $W_0$ and are processed by the receiver's carrier/channel tracking and verification rules (Section~\ref{sec:receiver}).
\item \textbf{Pilot Insertion Rule}: After the preamble, the $N_{\text{data}}$ data symbols are divided into $N_{\text{blk}}=16$ blocks of 4 symbols each. A pilot symbol is inserted at the beginning of each block. The structure of a single block is: \texttt{Pilot + Data + Data + Data + Data}.
\item \textbf{Number of Spread Symbols per Frame}: $N_{\text{sym}}=N_{\text{pre}}+N_{\text{pilot}}+N_{\text{data}}=82$.
\item $BW_{\text{hop}}$ (Hopping Bandwidth): $8\,\text{kHz}$. The hop range is centered at baseband $0\,\text{Hz}$ and spans $\pm 4\,\text{kHz}$.
\item $N_{\text{bins}}$ (Frequency Bins): $256$. Frequency resolution is $\Delta f = BW_{\text{hop}}/N_{\text{bins}} \approx 31.25\,\text{Hz}$.
\item \texttt{Domain} (AES-CTR Domain Separation): Fixed 32-bit constant \texttt{0x424C5443} (ASCII ``BLTC'').
\item $L_{\text{sym}}$ (Spread Symbol Length): Fixed at $L_{\text{sym}}=\mathit{SF}=1024$ chips for the information-bearing segment of each spread symbol.
\item $L_{\text{jitter}}$ (Guard Interval / Jitter Range): A per-symbol guard/noise segment appended after each spread symbol, with
  \[
    j_\ell \in [40,90]\ \text{chips},\quad \ell=0,\dots,N_{\text{sym}}-1.
  \]
  This corresponds to $8\,\text{ms}\sim 18\,\text{ms}$ at $R_c=5\,\text{kcps}$.
\item $T_{\text{sym,avg}}$ (Average Symbol Duration): Approximately $1024+65=1089$ chips, i.e.\ $\approx 217.8\,\text{ms}$.
\item $L_{\text{frame}}$ (Total Frame Length): Due to the per-symbol guard/noise insertion, the total frame length in chips is no longer a fixed constant. Instead,
  \[
    L_{\text{frame}}=\sum_{\ell=0}^{N_{\text{sym}}-1} (L_{\text{sym}} + j_\ell),
  \]
  while the symbol count remains fixed at $N_{\text{sym}}=82$. The number of samples is $L_{\text{frame}}\cdot \mathit{OSF}$ (plus optional tail padding, if present).
\item \textbf{Walsh Code Definition}: Uses the row vectors $W_m[j]\in\{+1,-1\}$ of the order-$\mathit{SF}=1024$ Sylvester-Hadamard matrix, where
  \[
    W_m[j]=(-1)^{\mathrm{popcount}(m \mathbin{\&} j)\bmod 2},\quad m\in[0,1023],\; j\in[0,1023].
  \]
  This protocol only uses the first $M_W$ orthogonal rows for $m\in[0,255]$. $W_0[j]\equiv +1$.
\item \textbf{RRC Pulse Shaping}:
  \begin{itemize}
  \item Roll-off: $\alpha=0.25$.
  \item Span: $S_{\text{rrc}}=6$ symbols.
  \item Taps: $L_{\text{rrc}} = S_{\text{rrc}}\cdot \mathit{OSF}+1 = 31$.
  \item Group delay (samples): $D_{\text{rrc}}=(L_{\text{rrc}}-1)/2=15$.
  \end{itemize}
\item \textbf{Tail and Ramp-Down}:
  \begin{itemize}
  \item Tail padding: $N_{\text{tail}}=8$ symbols (each symbol is $\mathit{SF}$ chips).
  \item Ramp-down duration: $T_{\text{ramp}}=20\,\text{ms}$.
  \end{itemize}
\item \textbf{Guard/Noise Filling}:
  \begin{itemize}
  \item Guard/noise chip power relative to data chips: $-3\,\text{dB}$.
  \item Guard soft-ramp length: $L_{\text{GI,ramp}}=8$ chips (capped by $j_\ell/2$ per symbol).
  \end{itemize}
\item \textbf{Receiver Acquisition Parameters}:
  \begin{itemize}
  \item Acquisition FFT size: $N_{\text{FFT}}=32768$.
  \item CFO search half-bandwidth: $f_{\text{search}}=8000\,\text{Hz}$.
  \item Hybrid verification threshold multiplier: $\gamma_{\text{hyb}}=10$.
  \item Intra-epoch start-offset stride: $N_{\text{off,stride}}=4$ samples.
  \item Fine-CFO grid: $\Delta f_{\text{fine}}=0.25\,\text{Hz}$ with $q\in\{-8,\dots,8\}$.
  \item Timing-refinement half-window: $T_{\text{rake}}=4\,\text{ms}$ (equivalently $N_{\text{rake}}=\lfloor F_s T_{\text{rake}}\rfloor$ samples).
  \item Pilot timing verification window: $N_{\text{pilot,win}}=32$ samples.
  \item Requested RAKE-finger output count: $N_{\text{finger,req}}$ with
    \[
      N_{\text{finger,out}}=\min(\max(N_{\text{finger,req}},1),16).
    \]
  \end{itemize}
\item \textbf{Receiver Tracking Parameters}:
  \begin{itemize}
  \item Numerical regularization: $\epsilon=10^{-18}$.
  \item Post-acquisition frame-start refinement window: $N_{\Delta n}=6$ samples (i.e.\ $\Delta n\in[-6,6]$).
  \item Maximum number of tracking/combining fingers: $N_{\text{finger,max}}=3$.
  \item Carrier-tracking frequency bank: $\Delta f_h \in \{-4.0,-3.75,\dots,+4.0\}\,\text{Hz}$ with step $0.25\,\text{Hz}$.
  \item Carrier-tracking frequency limit: $f_{\text{PLL,max}}=200\,\text{Hz}$.
  \item Frequency-snap rule constants: $\mathrm{conf}_{\text{freq,min}}=0.15$, $f_{\text{snap,tol}}=0.125\,\text{Hz}$, $N_{\text{snap}}=3$, $f_{\text{snap,min}}=0.75\,\text{Hz}$.
  \item Pilot cycle-slip magnitude threshold factor: $\eta_{\text{cs}}=0.25$ (threshold $|z_\ell|>\eta_{\text{cs}}M_{\text{pre}}$ on pilots).
  \item Pilot-aided channel-gain smoothing factor: $\alpha_A=\frac{1}{16}$.
  \item DLL early/late spacing: $0.5$ chips. DLL step limit: $\pm 2000\,\mathrm{ppm}$ around the nominal sample step.
  \end{itemize}
\end{itemize}

\section{Core Primitive: Cryptographically Secure Spreading Sequence Generator (CSPRNG Spreading)}

AES-CTR mode is used as the Pseudo-Random Number Generator (PRNG).

\subsection{Generation Function \texttt{GenCode(Key, TimeIndex, N\_sym)}}

\textbf{Inputs}:
\begin{itemize}
\item $\mathit{Key}$: The shared secret key.
\item $\mathit{TimeIndex}$: An integer time counter for the transmission instant.
\item $N_{\text{sym}}$: The number of spread symbols per frame (fixed at 82).
\end{itemize}

\textbf{Outputs}:
\begin{enumerate}
\item Cryptographic chip masks $C_{\text{seq}}$ for de-spreading, organized as $N_{\text{sym}}$ blocks of length $\mathit{SF}=1024$ chips.
\item Frequency control sequence $F_{\text{seq}}=\{f_0,f_1,\dots,f_{N_{\text{sym}}-1}\}$ (Hz).
\item Jitter/guard control sequence $J_{\text{seq}}=\{j_0,j_1,\dots,j_{N_{\text{sym}}-1}\}$ (chips).
\end{enumerate}

\textbf{Process}:
\begin{enumerate}
\item \textbf{Construct Nonce/Counter}: Follows the standard CTR ``Nonce + Counter'' structure.
  \begin{itemize}
  \item The 96-bit Nonce is defined as $\mathit{Nonce} = (\mathit{TimeIndex} \parallel \mathit{Domain})$, where $\mathit{TimeIndex}$ is 64-bit big-endian and $\mathit{Domain}$ is a 32-bit big-endian constant \texttt{0x424C5443} (ASCII for ``BLTC'').
  \item A 32-bit BlockCounter starts from 0 and increments for each block, forming the AES block input $\mathit{IV}_{\text{block}} = (\mathit{Nonce} \parallel \mathit{BlockCounter})$.
  \end{itemize}
\item \textbf{AES-CTR Generation}: Encrypt an all-zero data stream to produce the keystream bytes.
\item \textbf{Structured Consumption (Per Symbol)}: For the $\ell$-th spread symbol ($\ell=0,\dots,N_{\text{sym}}-1$), consume the keystream as:
  \begin{enumerate}
  \item \textbf{Data Block}: Take the next $1024$ bits and map them (MSB-first per byte) to the bipolar chip mask $C_{\text{seq}}[\ell][j]\in\{+1,-1\}$ using \texttt{0 -> +1}, \texttt{1 -> -1}.
  \item \textbf{Control Block}: Take the next $16$ bits as metadata:
    \begin{itemize}
    \item High 8 bits: $V_{\text{freq}}\in[0,255]$
    \item Low 8 bits: $V_{\text{jitter}}\in[0,255]$
    \end{itemize}
  \item \textbf{Parameter Mapping}:
    \begin{itemize}
    \item Hop offset:
      \[
        f_\ell=\left(\frac{V_{\text{freq}}}{256.0}-0.5\right)\cdot BW_{\text{hop}}\quad (\text{Hz}).
      \]
    \item Guard length:
      \[
        j_\ell =
        \begin{cases}
        40 + \left\lfloor \dfrac{V_{\text{jitter}}\cdot 50}{256}\right\rfloor & V_{\text{jitter}}\in[0,254] \\
        40 + 50 & V_{\text{jitter}}=255
        \end{cases}
        \quad (\text{chips}).
      \]
      Therefore $j_\ell\in[40,90]$.
    \end{itemize}
  \end{enumerate}
\item \textbf{Endianness and Bit Order Convention}:
  \begin{itemize}
  \item \texttt{TimeIndex} is encoded as \textbf{64-bit big-endian}.
  \item \texttt{Domain} and \texttt{BlockCounter} are encoded as \textbf{32-bit big-endian}.
  \item The AES-CTR output bytes are expanded into a bitstream with \texttt{bitorder=big} (MSB-first for each byte).
  \end{itemize}
\item \textbf{Non-Reuse Constraint}: Under a fixed shared key $K_{\text{sec}}$, each transmission uses a unique $\mathit{TimeIndex}$ and the transmitter does not transmit twice with the same $\mathit{TimeIndex}$. With $\mathit{TimeIndex}=\left\lfloor t_{\text{tx}}/\mathit{IV}_{\text{res}}\right\rfloor$, consecutive transmissions are separated by at least $\mathit{IV}_{\text{res}}$ in $t_{\text{tx}}$ to ensure distinct $\mathit{TimeIndex}$ values.
\end{enumerate}

\section{Transmitter Protocol}

Each transmitted frame is parameterized by a TimeIndex $\mathit{TI}_{\text{tx}}=\left\lfloor t_{\text{tx}}/\mathit{IV}_{\text{res}}\right\rfloor$, where $t_{\text{tx}}$ is the transmitter local time at the start of the frame. Successful reception requires $\mathit{TI}_{\text{tx}}\in\mathcal{T}$ (Section~\ref{sec:receiver}).

\subsection{Step A0: Frame Format}

This protocol uses a fixed \textbf{symbol count} per frame ($N_{\text{sym}}=82$). Due to per-symbol guard/noise insertion, the \textbf{total sample length} of a frame is variable and is determined by the generated $J_{\text{seq}}$ for the frame's $\mathit{TimeIndex}$.

\begin{itemize}
\item \textbf{Information Bits (uncoded)}: $U$, with a length of $K=256$ bits, composed of the following parts:
  \begin{enumerate}
  \item \textbf{Header (plaintext/structural fields)}:
    \begin{itemize}
    \item \texttt{Ver}: Protocol Version (4 bits).
    \item \texttt{Type}: Message Type (4 bits).
    \item \texttt{Len}: Payload Length (8 bits, in Bytes).
    \end{itemize}
  \item \textbf{Payload}: Service data bits.
  \item \textbf{CRC}: A CRC-32C is computed over the Header+Payload for post-decoding verification.
  \item \textbf{Padding}: If Header+Payload+CRC does not fill up $K$ bits, it is padded with zeros.
  \item \textbf{Byte/Bit Order}:
    \begin{itemize}
    \item $U$ is constructed as a 32-byte array: \texttt{Header(2) || Payload(Len) || CRC32C(4) || Pad(0)} to 32 bytes.
    \item Bytes are expanded MSB-first into bits.
    \item CRC32C is appended as 4 bytes in big-endian byte order.
    \end{itemize}
  \end{enumerate}

\item \textbf{Polar Encoding}: $U$ is encoded with a fixed code rate $R_{\text{FEC}}=1/2$ into a codeword $B$ of length $M=512$ (coded bits) using a fixed Polar code $(N,K)=(512,256)$.

  \textbf{Frozen Set and Bit Mapping}:
  \begin{enumerate}
  \item Construct a length-$N$ vector $\mathbf{u}=\{u_0,\dots,u_{N-1}\}$ with a fixed frozen set $\mathcal{F}$ of size $N-K$. For all $i\in\mathcal{F}$, set $u_i=0$. For all $i\notin\mathcal{F}$ (the information set), fill $u_i$ with the next bit from $U$ in increasing index order.
  \item The frozen set $\mathcal{F}$ is determined using Polarization Weight (PW) ranking for $N=512$ with $\beta=2^{1/4}$:
    \[
      w(i) = \sum_{j=0}^{\log_2 N - 1} b_j(i)\,\beta^j,\quad b_j(i)\in\{0,1\} \text{ is bit } j \text{ of } i.
    \]
    The $K$ indices with the largest $w(i)$ form the information set; all other indices are frozen.
  \item The codeword $\mathbf{B}=\{b_0,\dots,b_{N-1}\}$ is defined by the standard Arikan Polar transform applied to $\mathbf{u}$.
  \end{enumerate}

\item \textbf{Full-Frame Interleaving}: The interleaved sequence $B^{(\pi)}=\{b^{(\pi)}_0,\dots,b^{(\pi)}_{M-1}\}$ is defined by the permutation $\pi(\cdot)$ over the full codeword:
  \[
    b^{(\pi)}_j = b_{\pi(j)},\quad \pi(j) = (A\cdot j + B)\bmod M,\quad A=109,\; B=37.
  \]
  For $M=2^9$ with odd $A$, $\pi(\cdot)$ is a bijection.

\item \textbf{Symbol Mapping (256-ary Walsh)}: $B^{(\pi)}$ is grouped into $N_{\text{data}}=64$ symbol indices $m_q\in[0,255]$ of $k=8$ bits each (MSB first):
  \[
    m_q=\sum_{t=0}^{k-1} b^{(\pi)}_{qk+t}\cdot 2^{k-1-t},\quad q=0,\dots,N_{\text{data}}-1.
  \]
  These $N_{\text{data}}$ indices are mapped to $N_{\text{data}}$ orthogonal Walsh spread symbols.

\item \textbf{Symbol Arrangement (Preamble + Pilots + Data)}:
  \begin{itemize}
  \item Spread symbols are numbered $\ell=0,\dots,N_{\text{sym}}-1$.
  \item $\ell=0$ and $\ell=1$ are the two Preamble symbols.
  \item Following the preamble are $N_{\text{blk}}=16$ blocks, numbered $r=0,\dots,15$. Each block consists of 5 spread symbols: 1 pilot + 4 data.
  \item The position of the pilot symbol for block $r$ is
    \[
      \ell_{\text{pilot}}(r)=2+5r.
    \]
  \item The position of the $q$-th data symbol ($q=0,\dots,63$) is determined as follows: let $r=\lfloor q/4\rfloor$ and $s=q\bmod 4$, then
    \[
      \ell_{\text{data}}(q)=2+5r+1+s.
    \]
  \end{itemize}
\end{itemize}

\subsection{Step A: Payload Preparation}

The mapping from (\texttt{Header}, \texttt{Payload}) to $\{m_q\}_{q=0}^{N_{\text{data}}-1}$ is defined by Step~A0.

\subsection{Step B: Determine Transmission Time Seed}

The TimeIndex is defined by
\[
  \mathit{TI}_{\text{tx}} = \left\lfloor \frac{t_{\text{tx}}}{\mathit{IV}_{\text{res}}} \right\rfloor.
\]

\subsection{Step C: Generate Full-Frame Spreading Code}

Call the generation function:
\[
  (C_{\text{seq}}, F_{\text{seq}}, J_{\text{seq}}) = \texttt{GenCode}(K_{\text{sec}}, \mathit{TI}_{\text{tx}}, N_{\text{sym}}).
\]
where:
\begin{itemize}
\item $C_{\text{seq}}$ provides the $N_{\text{sym}}$ per-symbol chip masks (each of length $\mathit{SF}=1024$).
\item $F_{\text{seq}}$ provides the per-symbol hop offsets $f_\ell$ (Hz).
\item $J_{\text{seq}}$ provides the per-symbol guard lengths $j_\ell$ (chips).
\end{itemize}

\subsection{Step D: Walsh Orthogonal Spreading Modulation}

The cryptographic chip sequence is used as a ``mask/scrambling code'' and is combined with the orthogonal Walsh rows to generate the transmitted chips for each spread symbol.

The chips for the $\ell$-th spread symbol ($\ell=0,\dots,N_{\text{sym}}-1$) are denoted as $S[\ell\cdot \mathit{SF} + j]$, where $0\le j<\mathit{SF}$.

\begin{itemize}
\item \textbf{Preamble Symbols} ($\ell=0,1$): Fixed to use a Barker-2 synchronization word with Walsh index 0, i.e., $[+W_0,\,-W_0]$:
  \begin{itemize}
  \item \textbf{Preamble symbol 0} ($\ell=0$, $+W_0$):
    \[
      S[0\cdot \mathit{SF} + j] = C_{\text{seq}}[0\cdot \mathit{SF} + j].
    \]
  \item \textbf{Preamble symbol 1} ($\ell=1$, $-W_0$):
    \[
      S[1\cdot \mathit{SF} + j] = -C_{\text{seq}}[1\cdot \mathit{SF} + j].
    \]
  \end{itemize}
  These two preamble symbols do not carry information bits.

\item \textbf{Pilot Symbols} ($\ell=\ell_{\text{pilot}}(r)$, $r=0,\dots,15$): Fixed to use $W_0[j]\equiv +1$:
  \[
    S[\ell\cdot \mathit{SF} + j] = C_{\text{seq}}[\ell\cdot \mathit{SF} + j].
  \]
  Pilot symbols do not carry information bits.

\item \textbf{Data Symbols} ($\ell=\ell_{\text{data}}(q)$, $q=0,\dots,N_{\text{data}}-1$): Uses the $m_q$-th Walsh row. Let $\ell=\ell_{\text{data}}(q)$:
  \[
    S[\ell\cdot \mathit{SF} + j] = W_{m_q}[j]\cdot C_{\text{seq}}[\ell\cdot \mathit{SF} + j].
  \]
  This defines the information-bearing chip segment of each spread symbol (length $\mathit{SF}=1024$ chips). The transmitted chip stream is formed by appending a per-symbol guard/noise segment after each spread symbol (Step D1.5), resulting in a variable-length chip stream.
\end{itemize}

\subsection{Step D1.5: Guard Interval Insertion and Noise Filling}

For each spread symbol $\ell$, define the transmitted chip segment as the concatenation of the information-bearing chips and a guard/noise segment:
\[
  \{S[\ell\cdot \mathit{SF}+j]\}_{j=0}^{\mathit{SF}-1}\ \parallel\ \{G_\ell[k]\}_{k=0}^{j_\ell-1}.
\]

The guard/noise segment is defined as a deterministic real-valued sequence with variance $\sigma^2=10^{-3/10}$ (i.e.\ $\sigma=\sqrt{10^{-3/10}}$). Let $z_{\ell}[k]\sim\mathcal{N}(0,1)$ denote a standard-normal sequence defined by the PRNG mapping below, and define a soft-ramp window $w_{\text{GI},\ell}[k]\in[0,1]$ by:
\[
  r_\ell=\min\!\left(L_{\text{GI,ramp}},\left\lfloor \frac{j_\ell}{2}\right\rfloor\right),
\]
\[
  w_{\text{ramp}}[k]=\frac{1}{2}\left(1-\cos\left(\pi\frac{k}{r_\ell-1}\right)\right),\quad k=0,\dots,r_\ell-1,
\]
and
\[
  w_{\text{GI},\ell}[k]=
  \begin{cases}
  1 & r_\ell<2\\
  w_{\text{ramp}}[k] & r_\ell\ge 2\ \wedge\ 0\le k<r_\ell\\
  1 & r_\ell\ge 2\ \wedge\ r_\ell\le k<j_\ell-r_\ell\\
  w_{\text{ramp}}[j_\ell-1-k] & r_\ell\ge 2\ \wedge\ j_\ell-r_\ell\le k<j_\ell
  \end{cases}.
\]
The guard/noise chips are
\[
  G_\ell[k]=\sigma\cdot w_{\text{GI},\ell}[k]\cdot z_\ell[k],\quad k=0,\dots,j_\ell-1.
\]

\paragraph{Deterministic Normal Sequence Definition}
Define the 64-bit seed
\[
  s_{\ell,0}=\left(\mathit{TI}_{\text{tx}} \oplus (\ell\ll 32) \oplus \texttt{0xA7F031D359B29C11}\right)\oplus \texttt{0x9E3779B97F4A7C15}.
\]
Define the xorshift64* state recursion on 64-bit words (all operations modulo $2^{64}$):
\[
  x_{\ell,n,0}=s_{\ell,n},\quad x_{\ell,n,1}=x_{\ell,n,0}\oplus (x_{\ell,n,0}\gg 12),\quad x_{\ell,n,2}=x_{\ell,n,1}\oplus (x_{\ell,n,1}\ll 25),
\]
\[
  x_{\ell,n,3}=x_{\ell,n,2}\oplus (x_{\ell,n,2}\gg 27),\quad s_{\ell,n+1}=x_{\ell,n,3},\quad u_{64,\ell}[n]=x_{\ell,n,3}\cdot \texttt{0x2545F4914F6CDD1D}.
\]
Define the uniform sequence $u_\ell[n]\in(0,1)$ by
\[
  u_\ell[n]=\frac{(u_{64,\ell}[n]\gg 40)+1}{2^{24}+2}.
\]
Define the standard-normal sequence $z_\ell[k]$ by pairing uniforms $(u_\ell[2m],u_\ell[2m+1])$:
\[
  z_\ell[2m]=\sqrt{-2\ln(\max(u_\ell[2m],10^{-12}))}\cos(2\pi u_\ell[2m+1]),
\]
\[
  z_\ell[2m+1]=\sqrt{-2\ln(\max(u_\ell[2m],10^{-12}))}\sin(2\pi u_\ell[2m+1]).
\]

\subsection{Step D2: Pulse Shaping and Oversampling}

The baseband chip sequence is pulse-shaped.
\begin{itemize}
\item \textbf{Shaping Filter}: A Root-Raised-Cosine (RRC) filter $g_{\text{rrc}}(t)$ is used.
\item \textbf{Oversampling}: The shaped waveform is output at a sampling rate of $F_s = \mathit{OSF} \cdot R_c$.
\end{itemize}

Let $S_{\text{base}}[n]$ be the RRC-shaped complex baseband waveform at the sampling rate, generated from the chip stream that includes both the data segments and the guard/noise segments. Its length is $L_{\text{frame}}\cdot \mathit{OSF}$ (plus optional tail padding, if present).

RRC pulse shaping is applied before digital frequency rotation (Step D3). The filter taps are generated by the standard continuous-time RRC closed form sampled at $1/F_s$ and normalized to unit energy:

\begin{itemize}
\item Let $a=\alpha$ and define $t$ in chip periods. For each tap index $i=0,\dots,L_{\text{rrc}}-1$, define
  \[
    t_i = \frac{i - (L_{\text{rrc}}-1)/2}{\mathit{OSF}}.
  \]
\item The unnormalized impulse response is
  \[
    h(t)=\frac{\sin(\pi t(1-a)) + 4at\cos(\pi t(1+a))}{\pi t\left(1-(4at)^2\right)}.
  \]
  Special cases:
  \begin{itemize}
  \item At $t=0$:
    \[
      h(0) = 1-a+\frac{4a}{\pi}.
    \]
  \item At $t=\pm \frac{1}{4a}$:
    \[
      h\!\left(\pm\frac{1}{4a}\right)=\frac{a}{\sqrt{2}}\left[\left(1+\frac{2}{\pi}\right)\sin\!\left(\frac{\pi}{4a}\right)+\left(1-\frac{2}{\pi}\right)\cos\!\left(\frac{\pi}{4a}\right)\right].
    \]
  \end{itemize}
\item Define the sampled taps $\tilde h[i]=h(t_i)$ and the unit-energy normalized taps
  \[
    h[i]=\frac{\tilde h[i]}{\sqrt{\sum_{k=0}^{L_{\text{rrc}}-1} \tilde h[k]^2}}.
  \]
\end{itemize}

Chip upsampling is performed by placing each real chip at sample index $n = c\cdot \mathit{OSF}$ and inserting zeros at the intermediate samples. The pulse-shaped sequence is produced by causal FIR convolution
\[
  y[n]=\sum_{k=0}^{L_{\text{rrc}}-1} h[k]\cdot x[n-k].
\]

\subsection{Step D3: Digital Frequency Rotation and Phase Continuity}

Hopping is performed by complex rotation at baseband with a cumulative phase.

Let $f_\ell$ be the hop offset for symbol $\ell$ (from $F_{\text{seq}}$). Define a cumulative phase $\Phi[n]$ and a per-sample phase increment:
\[
  \Delta\omega_\ell = 2\pi\frac{f_\ell}{F_s}.
\]
For samples $n$ that fall within the $\ell$-th symbol's \textbf{time span} (including its guard/noise segment), update:
\[
  \Phi[-1]=0.
\]
\[
  \Phi[n] = \Phi[n-1] + \Delta\omega_\ell.
\]
Transmit waveform:
\[
  S_{\text{tx}}[n] = S_{\text{base}}[n]\cdot e^{j\Phi[n]}.
\]
\paragraph{Phase Continuity}
The phase process $\Phi[n]$ is continuous across symbol boundaries.

\subsection{Step E: Transmission}

This specification defines the complex baseband waveform $S_{\text{tx}}[n]$ at sampling rate $F_s$. Any RF upconversion and emission are outside the scope of this document.

\subsection{Step E2: Soft Shutdown}

\begin{enumerate}
\item \textbf{Tail Padding}: After transmitting the $N_{\text{sym}}$ symbols, append an additional $N_{\text{tail}}=8$ all-zero symbols (Zero Padding).
\item \textbf{Power Ramp-Down}: Let $L_{\text{ramp}}=\mathrm{round}(F_s\cdot T_{\text{ramp}})$ samples. Multiply the last $L_{\text{ramp}}$ samples by the half-cosine ramp
  \[
    w_{\text{tx}}[i]=\frac{1}{2}\left(1+\cos\left(\pi\frac{i}{L_{\text{ramp}}-1}\right)\right),\quad i=0,\dots,L_{\text{ramp}}-1,
  \]
  and define the ramped transmit sequence by $S_{\text{tx}}^{(\text{ramp})}[n]=S_{\text{tx}}[n]$ for $n<N-L_{\text{ramp}}$ and $S_{\text{tx}}^{(\text{ramp})}[N-L_{\text{ramp}}+i]=S_{\text{tx}}[N-L_{\text{ramp}}+i]\cdot w_{\text{tx}}[i]$.
\end{enumerate}

\section{Receiver Protocol}\label{sec:receiver}

The receiver input is a continuous complex baseband sample sequence $x[n]$ at sampling rate $F_s$.

The TimeIndex hypothesis set is $\mathcal{T}=\{\mathit{TI}_{\min},\dots,\mathit{TI}_{\min}+N_{\mathit{TI}}-1\}$. If $\mathit{TI}_{\min}$ and $N_{\mathit{TI}}$ are derived from a local-time estimate $t_{\text{rx}}$ and an uncertainty bound $W$, then
\[
  \mathit{TI}_{\min}=\left\lfloor \frac{t_{\text{rx}}-W}{\mathit{IV}_{\text{res}}}\right\rfloor,\quad
  N_{\mathit{TI}}=\left\lfloor \frac{t_{\text{rx}}+W}{\mathit{IV}_{\text{res}}}\right\rfloor-\mathit{TI}_{\min}+1.
\]

\subsection{Step A: Input Samples}

All receiver processing is defined on subsequences of $x[n]$. Buffering and RF front-end implementation details are outside the scope of this document.

\subsection{Step B: Blind Acquisition}

Blind acquisition maps $x[n]$ and $\mathcal{T}$ to an estimate of the frame parameters
\[
  \left(\widehat{\mathit{TI}}_{\text{tx}},\ \hat n_{\text{coarse}},\ \widehat{\Delta f},\ \{\hat n_{\text{finger},i}\}_{i=0}^{N_{\text{finger,sel}}-1}\right),
\]
where $N_{\text{finger,sel}}\le N_{\text{finger,out}}$.

\subsubsection{Step B.1: Hypothesis Sets}

Define the intra-epoch sample period
\[
  N_{\text{IV,samp}}=\mathrm{round}(F_s\cdot \mathit{IV}_{\text{res}}).
\]
Define the coarse frame-start hypothesis set (samples)
\[
  \mathcal{N}_{\text{coarse}}=\{0, N_{\text{off,stride}}, 2N_{\text{off,stride}}, \dots\}\cap [0, N_{\text{IV,samp}}).
\]

\subsubsection{Step B.2: Coarse Preamble Metric and Coarse CFO Estimate}

For a hypothesis pair $(\mathit{TI},n_0)\in\mathcal{T}\times\mathcal{N}_{\text{coarse}}$, define the locally generated schedules
\[
  (C_{\text{seq}}, F_{\text{seq}}, J_{\text{seq}})=\texttt{GenCode}(K_{\text{sec}}, \mathit{TI}, N_{\text{sym}}),
\]
and the symbol start offsets (relative to the frame start, in samples)
\[
  n_\ell = \left(\sum_{k=0}^{\ell-1} (\mathit{SF}+j_k)\right)\cdot \mathit{OSF},\quad \ell=0,\dots,N_{\text{sym}}-1.
\]

Let $s_0[n]$ and $s_1[n]$ be the locally generated TX-shaped, hopped preamble reference waveforms for preamble symbols $\ell=0,1$ (Barker-2 sign pattern $+W_0,-W_0$), each restricted to its $1024$-chip data segment. Let $\Delta n_1=(\mathit{SF}+j_0)\cdot \mathit{OSF}$ be the sample separation between the two preamble symbol starts. Phase continuity across the preamble boundary is defined by including the accumulated hop phase at the start of $s_1[n]$.

Define the preamble observation window $x_{n_0}[n]=x[n_0+n]$ and the sparse matched window
\[
  y_{\mathit{TI},n_0}[n]=
  \begin{cases}
  x_{n_0}[n]\cdot s_0^*[n] & n\in[0, \mathit{SF}\cdot\mathit{OSF}) \\
  0 & n\in[\mathit{SF}\cdot\mathit{OSF}, (\mathit{SF}+j_0)\cdot\mathit{OSF}) \\
  x_{n_0}[n]\cdot s_1^*[n-\Delta n_1] & n\in[\Delta n_1, \Delta n_1+\mathit{SF}\cdot\mathit{OSF})
  \end{cases}.
\]

Let $\widetilde y_{\mathit{TI},n_0}[n]$ be the length-$N_{\text{FFT}}$ zero-padded sequence defined by $\widetilde y_{\mathit{TI},n_0}[n]=y_{\mathit{TI},n_0}[n]$ for valid $n$ in the window above and $\widetilde y_{\mathit{TI},n_0}[n]=0$ otherwise. Define
\[
  Y_{\mathit{TI},n_0}[k]=\mathrm{FFT}_{N_{\text{FFT}}}\{\widetilde y_{\mathit{TI},n_0}[n]\}[k].
\]

Define the bin limit
\[
  k_{\max}=\left\lfloor \frac{\min(f_{\text{search}},F_s/2)\,N_{\text{FFT}}}{F_s}\right\rfloor,\quad k_{\max}\in[1,N_{\text{FFT}}/2-1],
\]
and the selected coarse bin
\[
  k_0=\arg\max_{k\in\{0,1,\dots,k_{\max}\}\cup\{N_{\text{FFT}}-k_{\max},\dots,N_{\text{FFT}}-1\}} |Y_{\mathit{TI},n_0}[k]|^2.
\]
Let $b_0\in[-k_{\max},k_{\max}]$ be the signed bin index corresponding to $k_0$. For $|b_0|<k_{\max}$, define the quadratic interpolation
\[
  \delta = \frac{1}{2}\frac{P_{-1}-P_{+1}}{P_{-1}-2P_0+P_{+1}},\quad b=b_0+\delta,
\]
with $P_{-1}=|Y_{\mathit{TI},n_0}[b_0-1]|^2$, $P_0=|Y_{\mathit{TI},n_0}[b_0]|^2$, $P_{+1}=|Y_{\mathit{TI},n_0}[b_0+1]|^2$ (circular indexing). The coarse CFO estimate is
\[
  \widehat{\Delta f}_{\text{coarse}}(\mathit{TI},n_0)=b\cdot \frac{F_s}{N_{\text{FFT}}}.
\]

Define the fine CFO grid
\[
  \mathcal{F}_{\text{fine}}(\mathit{TI},n_0)=\left\{\widehat{\Delta f}_{\text{coarse}}(\mathit{TI},n_0)+q\Delta f_{\text{fine}}:\ q\in\{-8,\dots,8\}\right\}.
\]
The refined CFO estimate $\widehat{\Delta f}(\mathit{TI},n_0)$ and the refined preamble energy $V_{\text{pre}}(\mathit{TI},n_0)$ are defined by
\[
  \widehat{\Delta f}(\mathit{TI},n_0)=\arg\max_{\Delta f\in\mathcal{F}_{\text{fine}}(\mathit{TI},n_0)}\left|\sum_n y_{\mathit{TI},n_0}[n]\cdot e^{-j2\pi\Delta f\,n/F_s}\right|^2,
\]
\[
  V_{\text{pre}}(\mathit{TI},n_0)=\max_{\Delta f\in\mathcal{F}_{\text{fine}}(\mathit{TI},n_0)}\left|\sum_n y_{\mathit{TI},n_0}[n]\cdot e^{-j2\pi\Delta f\,n/F_s}\right|^2.
\]

\subsubsection{Step B.3: Pilot Verification and Candidate Selection}

Define a noise power estimate from the raw baseband samples
\[
  \widehat{\sigma^2} = \frac{\mathrm{median}(|x[n]|^2)}{\ln 2}.
\]

Pilot symbol indices are defined by $\ell_{\text{pilot}}(r)=2+5r$, $r=0,\dots,N_{\text{pilot}}-1$. For each pilot $r$, let $p_r[n]$ be the locally generated TX-shaped, hopped pilot reference waveform for the pilot's $1024$-chip data segment. Define the pilot verification energy
\[
  E_r(\mathit{TI},n_0)=\max_{\Delta n\in[-N_{\text{pilot,win}},N_{\text{pilot,win}}]}\left|\sum_{n=0}^{\mathit{SF}\cdot\mathit{OSF}-1} x\!\left[n_0+n_{\ell_{\text{pilot}}(r)}+\Delta n+n\right]\cdot p_r^*[n]\cdot e^{-j2\pi\widehat{\Delta f}(\mathit{TI},n_0)\,n/F_s}\right|^2,
\]
and
\[
  V_{\text{pilot}}(\mathit{TI},n_0)=\sum_{r=0}^{N_{\text{pilot}}-1}E_r(\mathit{TI},n_0).
\]

Define the hybrid acceptance threshold
\[
  E_{\text{tot}} = E_{\text{pre}} + N_{\text{pilot}}\,E_{\text{sym}},\quad
  \Gamma = \gamma_{\text{hyb}}\cdot \widehat{\sigma^2}\cdot E_{\text{tot}}.
\]
The reference energies are defined by
\[
  E_{\text{pre}}=\sum_{n=0}^{\mathit{SF}\cdot\mathit{OSF}-1}|s_0[n]|^2+\sum_{n=0}^{\mathit{SF}\cdot\mathit{OSF}-1}|s_1[n]|^2,\quad
  E_{\text{sym}}=\sum_{n=0}^{\mathit{SF}\cdot\mathit{OSF}-1}|p_0[n]|^2.
\]
Define the accepted set
\[
  \mathcal{H}_{\text{acc}}=\left\{(\mathit{TI},n_0)\in\mathcal{T}\times\mathcal{N}_{\text{coarse}}:\ V_{\text{pre}}(\mathit{TI},n_0)+V_{\text{pilot}}(\mathit{TI},n_0)>\Gamma\right\}.
\]
If $\mathcal{H}_{\text{acc}}=\emptyset$, acquisition fails. Otherwise,
\[
  (\widehat{\mathit{TI}}_{\text{tx}},\hat n_{\text{coarse}})=\arg\max_{(\mathit{TI},n_0)\in\mathcal{H}_{\text{acc}}}\left(V_{\text{pre}}(\mathit{TI},n_0)+V_{\text{pilot}}(\mathit{TI},n_0)\right),
\]
with tie-break by smaller $(\mathit{TI},n_0)$.

\subsubsection{Step B.4: Timing Refinement and RAKE Finger Set}

Define the timing-refinement window (samples)
\[
  \mathcal{W}=\left\{\hat n_{\text{coarse}}-N_{\text{rake}},\dots,\hat n_{\text{coarse}}+N_{\text{rake}}\right\}.
\]
For each $\hat n\in\mathcal{W}$, define the two coherent preamble correlations
\[
  z_0(\hat n)=\sum_{n=0}^{\mathit{SF}\cdot\mathit{OSF}-1} x[\hat n+n]\cdot s_0^*[n],
\]
\[
  z_1(\hat n)=\sum_{n=0}^{\mathit{SF}\cdot\mathit{OSF}-1} x[\hat n+\Delta n_1+n]\cdot s_1^*[n],
\]
and the noncoherent timing score
\[
  \Lambda(\hat n)=|z_0(\hat n)|^2+|z_1(\hat n)|^2.
\]
The refined main-path start is
\[
  \hat n_0=\arg\max_{\hat n\in\mathcal{W}}\Lambda(\hat n),
\]
with tie-break by smaller $\hat n$.

Define the local-maximum set
\[
  \mathcal{M}=\left\{\hat n\in\mathcal{W}:\ \Lambda(\hat n)\ge \Lambda(\hat n-1)\ \wedge\ \Lambda(\hat n)>\Lambda(\hat n+1)\right\},
\]
with boundary cases interpreted by restricting to valid indices in $\mathcal{W}$. Define $\mathcal{M}_+=\{\hat n\in\mathcal{M}:\hat n\ge \hat n_0\}$.

Define the (unordered) RAKE finger set as any maximizer of
\[
  \mathcal{F}^\star\in\arg\max_{\substack{\mathcal{F}\subseteq\mathcal{M}_+\\ |\mathcal{F}|\le N_{\text{finger,out}}\\ \hat n_0\in\mathcal{F}\\ \forall \hat n_i\neq \hat n_j\in\mathcal{F}:\ |\hat n_i-\hat n_j|\ge \mathit{OSF}}}\ \sum_{\hat n\in\mathcal{F}}\Lambda(\hat n),
\]
with tie-break by choosing the lexicographically smallest sequence obtained by ordering $\mathcal{F}$ by descending $\Lambda(\hat n)$ and then ascending $\hat n$.

Let $N_{\text{finger,sel}}=|\mathcal{F}^\star|$. The ordered RAKE finger list $\{\hat n_{\text{finger},i}\}_{i=0}^{N_{\text{finger,sel}}-1}$ is defined as the elements of $\mathcal{F}^\star$ ordered by descending $\Lambda(\hat n)$ (tie-break: smaller $\hat n$).

The acquisition verification statistic is
\[
  V_{\text{acq}}=V_{\text{pre}}(\widehat{\mathit{TI}}_{\text{tx}},\hat n_{\text{coarse}})+V_{\text{pilot}}(\widehat{\mathit{TI}}_{\text{tx}},\hat n_{\text{coarse}}).
\]

\subsection{Step C: Matched Filtering, RAKE, and Tracking}

This step defines the matched-filtered, de-hopped sequence and the carrier/code tracking state variables that parameterize demodulation.

\subsubsection{Step C.1: Frame-Start Refinement}

Define the candidate shift set $\Delta\mathcal{N}=\{-N_{\Delta n},\dots,+N_{\Delta n}\}$ and the preference rank function
\[
  \rho(0)=0,\quad \rho(-d)=2d-1,\quad \rho(+d)=2d,\quad d=1,\dots,N_{\Delta n}.
\]
For each $\Delta n\in\Delta\mathcal{N}$, define the candidate frame start $\hat n_{\text{frm}}(\Delta n)=\hat n_0+\Delta n$ and shifted finger starts $\hat n_{\text{finger},i}(\Delta n)=\hat n_{\text{finger},i}+\Delta n$.

Let $\hat U(\Delta n)$ be the decoded information bits produced by applying Steps~C--D with frame start $\hat n_{\text{frm}}(\Delta n)$ and finger starts $\{\hat n_{\text{finger},i}(\Delta n)\}$. Let $\mathrm{CRC}(\hat U(\Delta n))\in\{0,1\}$ be the CRC-32C verification result.

Define the CRC-valid set $\Delta\mathcal{N}_{\text{crc}}=\{\Delta n\in\Delta\mathcal{N}:\mathrm{CRC}(\hat U(\Delta n))=1\}$. If $\Delta\mathcal{N}_{\text{crc}}\neq\emptyset$, then
\[
  \Delta n^\star=\arg\min_{\Delta n\in\Delta\mathcal{N}_{\text{crc}}}\rho(\Delta n).
\]
Otherwise define the verification statistic
\[
  V(\Delta n)=V_{\text{pre}}(\Delta n)+V_{\text{pilot}}(\Delta n),
\]
with
\[
  V_{\text{pre}}(\Delta n)=\sum_i\left|\left(\sum_{j=0}^{\mathit{SF}-1}\mathbf{u}_{i,0}[j]\right)+\left(\sum_{j=0}^{\mathit{SF}-1}\mathbf{u}_{i,1}[j]\right)\right|^2,
\]
\[
  V_{\text{pilot}}(\Delta n)=\sum_{\ell\in\mathcal{P}}|z_\ell|^2,
\]
where $\mathbf{u}_{i,\ell}[j]$ is defined in Step~C.4 and $z_\ell$ is defined in Step~C.5. The selected shift is
\[
  \Delta n^\star=\arg\max_{\Delta n\in\Delta\mathcal{N}} V(\Delta n),
\]
with tie-break by smaller $\rho(\Delta n)$.

The refined frame start and finger starts are
\[
  \hat n_{\text{frm}}=\hat n_{\text{frm}}(\Delta n^\star),\quad \hat n_{\text{finger},i}^{\star}=\hat n_{\text{finger},i}(\Delta n^\star).
\]

\subsubsection{Step C.2: CFO Derotation, De-hopping, and Matched Filtering}

Define the CFO-derotated sequence
\[
  x_{\text{cfo}}[n]=x[n]\cdot e^{-j2\pi\widehat{\Delta f}\,n/F_s}.
\]

Define the full-frame schedules from $\widehat{\mathit{TI}}_{\text{tx}}$:
\[
  (C_{\text{seq}}, F_{\text{seq}}, J_{\text{seq}})=\texttt{GenCode}(K_{\text{sec}}, \widehat{\mathit{TI}}_{\text{tx}}, N_{\text{sym}}).
\]
Define the symbol start offsets (samples)
\[
  n_\ell = \left(\sum_{k=0}^{\ell-1} (\mathit{SF}+j_k)\right)\cdot \mathit{OSF},\quad \ell=0,\dots,N_{\text{sym}}-1.
\]

Define the phase-continuous de-hopping phasor on the frame-relative sample index $m=n-\hat n_{\text{frm}}$. Let $\ell(m)$ be the unique symbol index whose time span (data+guard) contains $m$; for $m<0$, define $\ell(m)=0$. Define
\[
  \Phi[-1]=0.
\]
\[
  \Phi[m]=\Phi[m-1]+2\pi\frac{f_{\ell(m)}}{F_s},\quad D[m]=e^{-j\Phi[m]}.
\]
The de-hopped sequence is
\[
  x_{\text{dh}}[n]=x_{\text{cfo}}[n]\cdot D[n-\hat n_{\text{frm}}].
\]

Define the matched-filter output (receive RRC, same taps as the transmitter)
\[
  y[n]=\sum_k h[k]\cdot x_{\text{dh}}[n-k].
\]
The RRC group delay is $D_{\text{rrc}}$ samples and the cascade group delay is $D_{\text{cas}}=2D_{\text{rrc}}$ samples.

\subsubsection{Step C.3: RAKE Geometry and Windowing}

Let $N_{\text{finger,sel}}$ be the number of acquired finger positions and define
\[
  N_{\text{finger}}=\min(N_{\text{finger,max}}, N_{\text{finger,sel}}).
\]
Define the per-finger integer delay (samples) relative to the main finger:
\[
  \Delta d_i = \hat n_{\text{finger},i}^{\star} - \hat n_{\text{finger},0}^{\star},\quad i=0,\dots,N_{\text{finger}}-1.
\]

Define the pilot symbol index set
\[
  \mathcal{P}=\{\ell_{\text{pilot}}(r): r=0,\dots,N_{\text{pilot}}-1\},\quad \ell_{\text{pilot}}(r)=2+5r.
\]

For each spread symbol $\ell$, the data segment length is $\mathit{SF}$ chips and the guard/noise segment length is $j_\ell$ chips. All correlation and tracking quantities are defined on the $1024$-chip data window only.

\subsubsection{Step C.4: Despreading and MRC Combining}

Define the common timing correction state $\tau_\ell$ (samples) with initialization $\tau_0=0$ and update $\tau_{\ell+1}=\tau_\ell+\Delta t_\ell$, where $\Delta t_\ell$ is defined by the DLL in Step~C.6.

Define the tracked symbol start time for finger $i$ and symbol $\ell$ (samples)
\[
  t_{i,\ell}=\hat n_{\text{finger},i}^{\star}+D_{\text{cas}}+n_\ell+\tau_\ell.
\]
Define the tracked samples-per-chip $\Delta_{\text{chip}}=\mathit{step}_\ell/\mathit{SF}$. For each chip index $j=0,\dots,\mathit{SF}-1$, define the fractional sample index
\[
  p=t_{i,\ell}+j\Delta_{\text{chip}},\quad n_0=\lfloor p\rfloor,\quad a=p-n_0\in[0,1),
\]
and the linearly interpolated matched-filter sample
\[
  \mathbf{y}_{i,\ell}[j]=(1-a)\,y[n_0]+a\,y[n_0+1].
\]

Define the Barker-2 sign factor
\[
  b_\ell=
  \begin{cases}
  -1 & \ell=1\\
  +1 & \text{otherwise}
  \end{cases}.
\]
Define the hop-delay phase compensation factor
\[
  c_{i,\ell}=\exp\!\left(+j2\pi f_\ell \frac{\Delta d_i}{F_s}\right).
\]
The despread, hop-compensated chip vector is defined by
\[
  \mathbf{u}_{i,\ell}[j]=b_\ell\cdot \mathbf{y}_{i,\ell}[j]\cdot C_{\text{seq}}[\ell\cdot \mathit{SF}+j]\cdot c_{i,\ell}.
\]

Define the polarity state $p_\ell\in\{+1,-1\}$ with initialization $p_0=+1$ (updated by the cycle-slip rule in Step~C.5). The MRC weights and the combined chip vector are defined by
\[
  w_{i,\ell}=\frac{(p_\ell A_i)^*}{\sum_k |A_k|^2},\quad \mathbf{u}_{\ell}[j]=\sum_i w_{i,\ell}\cdot \mathbf{u}_{i,\ell}[j].
\]
The per-finger channel gains $A_i$ are initialized from the two preamble symbols and updated on pilot symbols with an exponential smoother ($\alpha_A=\frac{1}{16}$). For pilot index $r$ with symbol $\ell_{\text{pilot}}(r)$,
\[
  \tilde A_i^{(r)}=\frac{1}{\mathit{SF}}\sum_{j=0}^{\mathit{SF}-1}\mathbf{u}_{i,\ell_{\text{pilot}}(r)}[j]\cdot r_{\ell_{\text{pilot}}(r),h_{\text{best}}}[j],\quad
  A_i^{(r)}=(1-\alpha_A)A_i^{(r-1)}+\alpha_A\tilde A_i^{(r)}.
\]

\subsubsection{Step C.5: Carrier Tracking (PLL/Costas with Frequency Bank)}

Define the phase wrapping function
\[
  \mathrm{wrap}_{(-\pi,\pi]}(x)=(x+\pi)\bmod (2\pi)-\pi.
\]

The carrier-tracking state is $(\theta_\ell,\omega_\ell)$, where $\theta_\ell$ is the residual carrier phase at the start of the data window and $\omega_\ell$ is the residual phase advance across the data window of duration $T_{\text{data}}=\mathit{SF}/R_c$.

Define the per-chip compensation phasor
\[
  r_{\ell,h}[j]=\exp\left(-j\left(\theta_\ell + \frac{\omega_{\ell,h}}{\mathit{SF}}j\right)\right),\quad j=0,\dots,\mathit{SF}-1,
\]
with frequency-bank hypotheses
\[
  \Delta\omega_h = 2\pi\,\Delta f_h\,T_{\text{data}},\quad
  \omega_{\ell,h}=\min\!\left(\max\!\left(\omega_\ell+\Delta\omega_h,-\omega_{\max}\right),\omega_{\max}\right),
\]
and $\omega_{\max}=2\pi f_{\text{PLL,max}}T_{\text{data}}$.

For preamble/pilot symbols ($W_0$), define the hypothesis score
\[
  z_{\ell,h}=\sum_{j=0}^{\mathit{SF}-1}\mathbf{u}_\ell[j]\cdot r_{\ell,h}[j],\quad
  h_{\text{best}}=\arg\max_h |z_{\ell,h}|^2.
\]
For data symbols, define
\[
  R_{\ell,h}[m]=\sum_{j=0}^{\mathit{SF}-1}\mathbf{u}_\ell[j]\cdot r_{\ell,h}[j]\cdot W_m[j],\quad m\in[0,255],
\]
and
\[
  (h_{\text{best}},\hat m_\ell)=\arg\max_{h,m\in[0,255]} |R_{\ell,h}[m]|^2.
\]

Let $P_{(1)}\ge P_{(2)}$ be the largest and second-largest values of $|R_{\ell,h_{\text{best}}}[m]|^2$ over $m\in[0,255]$, and let $Q_{(1)}\ge Q_{(2)}$ be the largest and second-largest values of the selected hypothesis score (pilot: $|z_{\ell,h}|^2$; data: $\max_m |R_{\ell,h}[m]|^2$). Define
\[
  \mathrm{conf}_{\text{code}}=\frac{P_{(1)}-P_{(2)}}{P_{(1)}+\epsilon},\quad
  \mathrm{conf}_{\text{freq}}=\frac{Q_{(1)}-Q_{(2)}}{Q_{(1)}+\epsilon}.
\]

Define the symbol span factor
\[
  s_\ell = 1 + \frac{j_\ell}{\mathit{SF}}.
\]
Define the decision statistic
\[
  z_\ell=
  \begin{cases}
  z_{\ell,h_{\text{best}}} & \ell\in\{0,1\}\cup\mathcal{P}\\
  R_{\ell,h_{\text{best}}}[\hat m_\ell] & \text{otherwise}
  \end{cases},
\quad
  \tilde z_\ell=
  \begin{cases}
  z_\ell & \Re\{z_\ell\}\ge 0\\
  -z_\ell & \Re\{z_\ell\}<0
  \end{cases}.
\]
Define the phase error
\[
  e_\ell=\operatorname{atan2}(\Im\{\tilde z_\ell\},\Re\{\tilde z_\ell\}+\epsilon).
\]
Define the decision-directed scale
\[
  s_{\text{dd}}=
  \begin{cases}
  1.0 & \ell\in\{0,1\}\cup\mathcal{P}\\
  0.5 & \text{otherwise}
  \end{cases}.
\]

Given PLL gains $(K_p,K_i)$, define the carrier state update
\[
  \omega_{\ell+1}^{(0)}=\min\!\left(\max\!\left(\omega_\ell + (K_i s_{\text{dd}})\,e_\ell,-\omega_{\max}\right),\omega_{\max}\right),
\]
\[
  \theta_{\ell+1}=\mathrm{wrap}_{(-\pi,\pi]}\!\left(\theta_\ell + \omega_{\ell,h_{\text{best}}} s_\ell + (K_p s_{\text{dd}})\,e_\ell\right).
\]

Define the pilot cycle-slip indicator using a preamble magnitude reference $M_{\text{pre}}$:
\[
  \chi_\ell=
  \begin{cases}
  1 & \ell\in\mathcal{P}\ \wedge\ \Re\{z_\ell\}<0\ \wedge\ |z_\ell|>\eta_{\text{cs}}\,M_{\text{pre}}\\
  0 & \text{otherwise}
  \end{cases}.
\]
The cycle-slip corrected carrier phase and polarity state are defined by
\[
  \theta_{\ell+1}=\mathrm{wrap}_{(-\pi,\pi]}\!\left(\theta_{\ell+1}+\chi_\ell\pi\right),\quad
  p_{\ell+1}=(-1)^{\chi_\ell}p_\ell.
\]

Define the frequency-snap state $(\nu_\ell,c_\ell)$ with initialization $(\nu_0,c_0)=(0,0)$, where $\nu_\ell$ is the held bank offset (Hz) and $c_\ell\in\{0,\dots,N_{\text{snap}}\}$ is a consecutive-count state. Let $\Delta f_{\ell}=\Delta f_{h_{\text{best}}}$. The state update is
\[
  (\nu'_{\ell+1},c'_{\ell+1})=
  \begin{cases}
  (\nu_\ell,\min(c_\ell+1,N_{\text{snap}})) & \mathrm{conf}_{\text{freq}}\ge \mathrm{conf}_{\text{freq,min}}\ \wedge\ \Delta f_{\ell}\neq 0\ \wedge\ |\Delta f_{\ell}-\nu_\ell|\le f_{\text{snap,tol}}\\
  (\Delta f_{\ell},0) & \text{otherwise}.
  \end{cases}
\]
Define the snap event indicator
\[
  \psi_{\ell+1}=\mathbf{1}\{c'_{\ell+1}=N_{\text{snap}}\}.
\]
The post-snap frequency state and stored snap state are defined by
\[
  \omega_{\ell+1}=\min\!\left(\max\!\left(\omega_{\ell+1}^{(0)}+\psi_{\ell+1}\cdot \mathbf{1}\{|\nu'_{\ell+1}|\ge f_{\text{snap,min}}\}\cdot 2\pi\nu'_{\ell+1}T_{\text{data}},-\omega_{\max}\right),\omega_{\max}\right),
\]
\[
  (\nu_{\ell+1},c_{\ell+1})=
  \begin{cases}
  (0,0) & \psi_{\ell+1}=1\\
  (\nu'_{\ell+1},c'_{\ell+1}) & \psi_{\ell+1}=0
  \end{cases}.
\]

Given loop bandwidth $B_L$ (Hz) and damping $\zeta$, define
\[
  T_{\text{upd}}=\frac{\mathit{SF}+j_{\text{avg}}}{R_c},\quad j_{\text{avg}}=40+\frac{50}{2},
\]
\[
  \omega_n=2\pi B_L,\quad r=\exp(-\zeta\omega_n T_{\text{upd}}),\quad \phi=\omega_n\sqrt{1-\zeta^2}\,T_{\text{upd}},
\]
\[
  K_p = 1-r^2,\quad K_i = 1+r^2-2r\cos\phi.
\]
Fixed PLL values: $B_L=1.0\,\text{Hz}$, $\zeta=0.707$.

\subsubsection{Step C.6: Code Timing Tracking (Early/Late DLL)}

The DLL timing state is $\mathit{step}_\ell$ (samples per $1024$-chip data window) and $\tau_\ell$ (common timing offset, samples). Define
\[
  \mathit{step}_{\text{nom}}=\mathit{SF}\cdot \mathit{OSF},\quad
  \mathit{step}_{\min}=\mathit{step}_{\text{nom}}(1-2000\cdot 10^{-6}),\quad
  \mathit{step}_{\max}=\mathit{step}_{\text{nom}}(1+2000\cdot 10^{-6}).
\]
Define the early/late spacing (samples)
\[
  \Delta_{\text{EL}} = 0.5\cdot \mathit{step}_\ell/\mathit{SF}.
\]
Define $\Delta_{\text{chip}}=\mathit{step}_\ell/\mathit{SF}$.

Define the early/late sampling instants
\[
  p_E=t_{i,\ell}+j\Delta_{\text{chip}}-\Delta_{\text{EL}},\quad
  p_L=t_{i,\ell}+j\Delta_{\text{chip}}+\Delta_{\text{EL}}.
\]
Define
\[
  n_E=\lfloor p_E\rfloor,\quad a_E=p_E-n_E\in[0,1),\quad \mathbf{y}_{i,\ell}^{(E)}[j]=(1-a_E)\,y[n_E]+a_E\,y[n_E+1],
\]
\[
  n_L=\lfloor p_L\rfloor,\quad a_L=p_L-n_L\in[0,1),\quad \mathbf{y}_{i,\ell}^{(L)}[j]=(1-a_L)\,y[n_L]+a_L\,y[n_L+1].
\]

Define the despread, hop-compensated early/late chip vectors
\[
  \mathbf{u}_{i,\ell}^{(E)}[j]=b_\ell\cdot \mathbf{y}_{i,\ell}^{(E)}[j]\cdot C_{\text{seq}}[\ell\cdot \mathit{SF}+j]\cdot c_{i,\ell},\quad
  \mathbf{u}_{i,\ell}^{(L)}[j]=b_\ell\cdot \mathbf{y}_{i,\ell}^{(L)}[j]\cdot C_{\text{seq}}[\ell\cdot \mathit{SF}+j]\cdot c_{i,\ell}.
\]
Define the combined early/late vectors using the same weights as the prompt branch:
\[
  \mathbf{u}_{\ell}^{(E)}[j]=\sum_i w_{i,\ell}\cdot \mathbf{u}_{i,\ell}^{(E)}[j],\quad
  \mathbf{u}_{\ell}^{(L)}[j]=\sum_i w_{i,\ell}\cdot \mathbf{u}_{i,\ell}^{(L)}[j].
\]

Define the carrier-compensated early/late correlator outputs
\[
  z_E=
  \begin{cases}
  \sum_{j=0}^{\mathit{SF}-1}\mathbf{u}_{\ell}^{(E)}[j]\cdot r_{\ell,h_{\text{best}}}[j] & \ell\in\{0,1\}\cup\mathcal{P}\\
  \sum_{j=0}^{\mathit{SF}-1}\mathbf{u}_{\ell}^{(E)}[j]\cdot r_{\ell,h_{\text{best}}}[j]\cdot W_{\hat m_\ell}[j] & \text{otherwise}
  \end{cases},
\]
\[
  z_L=
  \begin{cases}
  \sum_{j=0}^{\mathit{SF}-1}\mathbf{u}_{\ell}^{(L)}[j]\cdot r_{\ell,h_{\text{best}}}[j] & \ell\in\{0,1\}\cup\mathcal{P}\\
  \sum_{j=0}^{\mathit{SF}-1}\mathbf{u}_{\ell}^{(L)}[j]\cdot r_{\ell,h_{\text{best}}}[j]\cdot W_{\hat m_\ell}[j] & \text{otherwise}
  \end{cases}.
\]

Define
\[
  d_\ell=\frac{|z_E|-|z_L|}{|z_E|+|z_L|+\epsilon},\quad
  e_{\text{DLL},\ell}=d_\ell\cdot\frac{\Delta_{\text{EL}}}{2}.
\]

Given DLL gains $(K_p^{\text{DLL}},K_i^{\text{DLL}})$ and decision-directed scale
\[
  s_{\text{dd}}^{\text{DLL}}=
  \begin{cases}
  1.0 & \ell\in\{0,1\}\cup\mathcal{P}\\
  0.25 & \text{otherwise}
  \end{cases},
\]
define the timing updates
\[
  \mathit{step}_{\ell+1}=\min\!\left(\max\!\left(\mathit{step}_\ell - (K_i^{\text{DLL}} s_{\text{dd}}^{\text{DLL}})\,e_{\text{DLL},\ell},\mathit{step}_{\min}\right),\mathit{step}_{\max}\right),
\]
\[
  \Delta t_\ell = -(K_p^{\text{DLL}} s_{\text{dd}}^{\text{DLL}})\,e_{\text{DLL},\ell}.
\]

Fixed DLL values: $B_L=0.6\,\text{Hz}$, $\zeta=0.707$ with the same pole-mapping definitions as the PLL.

\subsection{Step D: Walsh Demodulation and Polar (CA-SCL) Decoding}

\subsubsection*{1. Pilot-aided Carrier Tracking}

Pilot symbols (known $W_0$) are the elements of $\mathcal{P}$ and define the known-reference decision statistic $z_\ell$ in Step~C.5 and the channel-gain update in Step~C.4.

\subsubsection*{2. Walsh Matched Filtering (FHT)}

For each data symbol $\ell=\ell_{\text{data}}(q)$ ($q=0,\dots,N_{\text{data}}-1$) and each symbol-rate frequency hypothesis $k$ (from the bank in Step~C.5), the $M_W$-way correlation is defined by
\[
  R_{\ell,k}[m]=\sum_{j=0}^{\mathit{SF}-1}\mathbf{u}_{\ell}[j]\cdot r_{\ell,k}[j]\cdot W_m[j],\quad m\in[0,255].
\]
The selected hypothesis is
\[
  k_{\text{best}}=\arg\max_k \max_{m\in[0,255]} |R_{\ell,k}[m]|^2,
\]
and $R_{\ell,k_{\text{best}}}[m]$ defines the soft-demapping input.

\subsubsection*{3. Soft Information (LLR) Generation}

The decision-directed common phase rotation is defined by
\[
  \widetilde R_{\ell}[m] = R_{\ell,k_{\text{best}}}[m]\cdot e^{-j\arg\!\left(R_{\ell,k_{\text{best}}}[\hat m_\ell]\right)}.
\]
Let the metric be $D_{\ell}[m]=\Re\{\widetilde R_{\ell}[m]\}$. If the decision-directed rotation is undefined (zero magnitude), use $D_{\ell}[m]=\Re\{R_{\ell,k_{\text{best}}}[m]\}$. For the $t$-th bit within each symbol ($t=0,\dots,k-1$), calculate the max-log LLR:
\[
  \mathit{LLR}_{qk+t}=\max_{m:\, ((m\gg(k-1-t))\mathbin{\&}1)=0} D_{\ell}[m] - \max_{m:\, ((m\gg(k-1-t))\mathbin{\&}1)=1} D_{\ell}[m].
\]
The $N_{\text{data}}\cdot k=M$ LLRs correspond to the interleaved coded bits $B^{(\pi)}$. The deinterleaved LLRs are defined by the inverse permutation:
\[
  \mathit{LLR}_i = \mathit{LLR}^{(\pi)}_{\pi^{-1}(i)}.
\]

\subsubsection*{4. CA-SCL Polar Decoding and Output}

Let $\hat U$ be the information-bit estimate produced by CRC-aided SCL (CA-SCL) Polar decoding of the deinterleaved LLRs. List selection is defined by CRC-32C (embedded inside $U$): among the list candidates, $\hat U$ is the best-metric candidate that satisfies CRC; if no list candidate satisfies CRC, $\hat U$ is the best-metric candidate.

The decode is accepted if $\mathrm{CRC32C}(\hat U)=0$ under the CRC definition in Step~A0. If accepted, the receiver output is the payload extracted from $\hat U$ using the \texttt{Len} field; otherwise the packet is invalid.

\end{document}
